\documentclass[10pt,twocolumn]{article}

\usepackage[margin=0.60in,columnsep=0.22in]{geometry}
\usepackage{amsmath,amssymb,booktabs,enumitem,microtype}
\usepackage[font=small,labelfont=bf]{caption}
\usepackage{xcolor}
\usepackage[colorlinks=true,allcolors=blue!55!black]{hyperref}

\setlength{\parindent}{0pt}
\setlength{\parskip}{3pt}
\setlength{\emergencystretch}{2em}
\setlist{nosep,leftmargin=*}
\newcommand{\D}{\Delta_{60}}
\newcommand{\BR}{\operatorname{BR}}

\title{\vspace{-1.2em}\textbf{Perfect Mode Hal: Aggressive, Change-Aware\\
Opponent Exploitation in Exact DTH}}
\author{James Cenawood}
\date{August 20, 2026}

\begin{document}
\maketitle
\vspace{-1.3em}

\begin{abstract}
Perfect Mode (PM) Hal is a causal meta-policy for repeated, pure Drop the
Handkerchief (DTH). One certified tablebase remains the sole authority for each
continuation-adjusted $60\times60$ matrix. Around it, PM combines Adaptive
Hal's Bayesian counts, Perfect Hal's pattern experts, online change-point and
outcome models, and a compatible Aggro Hal recurrent network. Fixed Share
learns role-specific source and controller weights. Four modes move from an
exact shield to hard forecast responses, while every selected policy receives
a fresh local worst-case-loss measurement. An earlier 24-identity development
run scored 44 wins in 48 games, but its protocol lacked durable registration
and its evaluator trusted self-reported risk. It has no promotion authority.
This note instead registers a 56-identity confirmation protocol with independent
matrix recomputation, a no-Aggro ablation, and family and seat gates. At the
protocol commit that confirmation has not been executed. PM is an aggressive
candidate, not human mind reading, dominance, or a whole-game safety theorem.
\end{abstract}

\section{Scope and design philosophy}

At public state $x$, the exact DTH agent supplies matrix $M(x)$, equilibrium
policies $(x^\star,y^\star)$, and a saddle-gap certificate at most $10^{-6}$.
All actions are literal seconds $1,\ldots,60$. PM never estimates continuation
values and never sees an unrevealed simultaneous action.

The intended character is ``overbearing after evidence, composed after
surprise.'' Perfect contributes hard exploitation; Adaptive contributes
uncertainty and a constrained response frontier; Aggro contributes recurrent
features and a learned residual policy. The story evidence ledger motivates
competing personae, opponent mental models, and abrupt memory discontinuities.
PM translates that fiction into testable engineering---model competition,
role separation, recency, and change recovery---rather than treating a person
as one fixed psychological type. The final narrative ideal, to improve and
adapt without looking back, becomes persistent memory with rapid de-weighting
and recovery after changed play.

\section{Causal opponent model}

For each opponent role separately, PM maintains six mandatory forecast
sources: uniform recovery, the opponent's exact equilibrium policy, a decayed
Dirichlet model, Perfect's online expert mixture, categorical Bayesian online
change-point detection (BOCPD), and an outcome-conditioned model. A seventh
source is Aggro's GRU forecast when a strict compatible checkpoint is present.
After action $a_t$ is revealed, source $k$ receives prequential log score and
Fixed Share updates
\begin{align*}
 \widetilde w_{k,t+1}&\propto w_{k,t}
       \exp\!\left(\eta\log q_{k,t}(a_t)\right),\\
 w_{k,t+1}&=(1-\alpha)\widetilde w_{k,t+1}
 +\frac{\alpha}{K-1}\left(1-\widetilde w_{k,t+1}\right),
\end{align*}
with $\alpha=.04$. Thus every source retains recovery mass when an opponent
switches. BOCPD carries a truncated Dirichlet--multinomial run-length posterior
with hazard $.08$. The fused forecast is
$\widehat q_t=\sum_k w_{k,t}q_{k,t}$.

Aggression confidence is the declared heuristic
\[
 c_t=c_{\max}\frac{n_{\rm eff}}{4+n_{\rm eff}}
 (1-P_t(\text{change}))(1-\mathrm{JS}_t)s_t,
\]
where $\mathrm{JS}_t$ is normalized source disagreement and
$s_t=\exp[-\max(0,\overline{\mathrm{NLL}}_t-\log60)]$. Only the evidence
fraction has a data-biased-response interpretation; the product is a project
controller, not a published theorem.

\section{Aggression with measured risk}

For Dropper, the certified $\epsilon$ candidate solves
\[
 \max_{x\in\D}x^{\mathsf T}M\widehat q_t
 \quad\text{s.t.}\quad
 (M^{\mathsf T}x)_j\ge L-\epsilon\quad\forall j,
\]
where $L=\min_j(M^{\mathsf T}x^\star)_j$ is freshly recomputed. Checker uses
the symmetric upper constraint. The bank also contains exact equilibrium,
hard responses to fused and Perfect forecasts, and Aggro's direct policy. For
a direct Dropper policy $x$, actual local loss is
$\max(0,L-\min_j(M^{\mathsf T}x)_j)$; every direct and frontier candidate is
independently checked against the current matrix.

One common bank of heuristic ensemble perturbations must give positive expected
improvement and at least $.55$ positive-draw support. These draws are not a
calibrated posterior, so the support fraction is not an improvement probability.
Cold start or a change shock enters \emph{shield}, where only the exact face is
admissible. The registered v3 caps are \emph{probe} $.05$, \emph{press} $.50$,
and \emph{dominate} $2.0$, inside a $12.0$ per-game charge budget. Controller
weights receive full-information counterfactual payoff feedback after reveal.
This is deliberately aggressive: the budget is an explicit risk appetite, not
a guarantee that losses add across a DTH game. Ganzfried--Sandholm's cumulative
ledger concerns repetitions of one stage game; DTH's matrices change with the
state and already contain overlapping continuation values.

\section{Development audit and registered confirmation}

The v1 and v2 files were created and run before this feature branch registered
them in Git. File status text and mutable modification times cannot establish
preregistration. Both runs are therefore development evidence. The v2 run used
24 parameter identities across six opponent families. Each identity produced a
common seeded seat pair for PM, Perfect, Adaptive, Aggro, and Exact; family plus
parameter seed was the bootstrap unit (10,000 replicates). Stopped games scored
$.5$.

\begin{table}[h]
\centering
\caption{V2 development outcomes (48 games per controller).}
\small
\begin{tabular}{lrrr}
\toprule
Controller & Score & W--L--S & Mean halves\\
\midrule
PM Hal     & .9167 & 44--4--0  & 6.63\\
Perfect    & .8958 & 43--5--0  & 5.56\\
Adaptive   & .6563 & 31--16--1 & 7.10\\
Aggro      & .6458 & 31--17--0 & 7.23\\
Exact      & .5417 & 26--22--0 & 9.46\\
\bottomrule
\end{tabular}
\end{table}

Paired PM-minus-baseline score differences were $+.3750$
$[.2083,.5417]$ versus Exact, $+.2604$ $[.1354,.3854]$ versus Adaptive,
$+.2708$ $[.1042,.4375]$ versus Aggro, and $+.0208$
$[-.0417,.0833]$ versus Perfect. Those intervals are descriptive. The run cannot
promote PM, and even its nominal gate failed the nonnegative lower-bound
condition against every component.

Across 318 PM decisions, shield/probe/press/dominate counts were
$56/46/154/62$. The controller reported maximum local loss $.2561$ and no
understated-loss or game-budget violations, but the v2 evaluator did not
independently reconstruct every selected policy's loss. PM's identity-clustered expected
NLL was $3.481$ $[2.904,3.978]$. Fused-minus-source NLL was $-.252$
$[-.341,-.162]$ against Perfect and $-.614$ $[-1.184,-.119]$ against uniform;
negative is better. The fused forecast beat every individual source on this
metric. Change detection remained weak: it detected 10 of 86 synthetic truth
shifts, with one false-alarm onset. Strong play did not validate that detector.

The v3 confirmation is registered in Git before execution. It binds the
controller, recurrent checkpoint, and tablebase by SHA-256; requires a clean
worktree; records the executing commit; and assigns four new seeds to all 14
registered synthetic families, for 56 independent identities and 112 games per
controller. It compares PM with Exact, Perfect, Adaptive, Aggro, and PM without
Aggro. Promotion requires zero independent risk-audit violations, nonnegative
clustered lower bounds against Exact, every component, and the no-Aggro
ablation, plus nonnegative point differences in both seats and every family.
The evaluator reloads every continuation-adjusted matrix and recomputes selected
worst-case loss, charge contracts, mode caps, cumulative charges, and per-game
budgets. No v3 result appears in this protocol-registration version of the note.

\section{Human-play interpretation and limits}

The psychology is operational: reinforcement, belief learning, response to
one's own revealed action, recency, periodicity, and change are competing
causal hypotheses. A human can counter-adapt, deceive, or change strategy; PM
must preserve recovery mass and retreat on surprise. It should never label a
participant's personality or claim privileged access to intent.

A human study must keep opponent identity across games, log the sequence
forecast before each reveal, score realized NLL/Brier without simulator-truth
distributions, and cluster outcomes by participant. Change delay and false
alarms, role slices, calibration, and game score are co-primary diagnostics.
Bounded recursive reasoning is one possible population archetype, not a label
the controller may assign permanently to a participant.
The present measured evidence is synthetic development evidence from 24
identities and known families. The recurrent checkpoint saw 2,653 training
decisions; its causal contribution remains an open question until the registered
no-Aggro ablation is read. PM is feature-complete for human play as an
aggressive candidate, while promotion awaits the registered synthetic result
and then participant-level human evidence.

\begingroup
\footnotesize
\setlength{\parskip}{0pt}
\begin{thebibliography}{9}
\bibitem{rnr} M. Johanson, M. Zinkevich, and M. Bowling.
\href{https://papers.nips.cc/paper_files/paper/2007/file/6e7b33fdea3adc80ebd648fffb665bb8-Paper.pdf}{Computing Robust Counter-Strategies}. NeurIPS, 2007.
\bibitem{dbr} M. Johanson and M. Bowling.
\href{https://proceedings.mlr.press/v5/johanson09a.html}{Data Biased Robust Counter Strategies}. AISTATS, 2009.
\bibitem{share} M. Herbster and M. Warmuth.
\href{https://doi.org/10.1023/A:1007424614876}{Tracking the Best Expert}. Machine Learning, 1998.
\bibitem{bocpd} R. Adams and D. MacKay.
\href{https://arxiv.org/abs/0710.3742}{Bayesian Online Changepoint Detection}, 2007.
\bibitem{ewa} C. Camerer and T. Ho.
\href{https://doi.org/10.1111/1468-0262.00054}{Experience-Weighted Attraction Learning}. Econometrica, 1999.
\bibitem{mentalize} A. Hampton, P. Bossaerts, and J. O'Doherty.
\href{https://doi.org/10.1073/pnas.0711099105}{Mentalizing-related computations in strategic interaction}. PNAS, 2008.
\bibitem{delta} M. Nassar et al.
\href{https://doi.org/10.1523/JNEUROSCI.0822-10.2010}{A Bayesian delta-rule model of changing environments}. J. Neuroscience, 2010.
\bibitem{ch} C. Camerer, T.-H. Ho, and J.-K. Chong.
\href{https://doi.org/10.1162/0033553041502225}{A Cognitive Hierarchy Model of Games}. QJE, 2004.
\bibitem{scores} T. Gneiting and A. Raftery.
\href{https://doi.org/10.1198/016214506000001437}{Strictly Proper Scoring Rules}. JASA, 2007.
\bibitem{safe} S. Ganzfried and T. Sandholm.
\href{https://doi.org/10.1145/2716322}{Safe Opponent Exploitation}. ACM TEAC, 2015.
\end{thebibliography}
\endgroup

\end{document}
